\section{The order restriction}\label{sec:reduction}

We begin with the classical restriction that a circulant Hadamard order
greater than one has the form \(4u^2\) with \(u\) odd. The descent below
is a group-ring version of Turyn's argument for the \(2\)-primary
coefficients \cite[Lemma~6 through Theorem~8 and its Corollary,
pp.~334--336]{Turyn1965}. We first isolate the cyclotomic divisibility
fact that the descent uses.

Let \(H\) be a real circulant Hadamard matrix of order \(n>1\), and
represent its first row by
\[
 h=\sum_{j=0}^{n-1}h_jX^j\in\Z[C_n].
\]
The coefficient of \(X^a\) in \(hh^*\) is
\(\sum_j h_jh_{j-a}\), the inner product of row zero and row \(a\).
Consequently orthogonality is exactly
\begin{equation}\label{eq:hadamard-norm}
 hh^*=n.
\end{equation}
Augmentation gives \(n=h(1)^2\), so \(n\) is a square. The inner
product of two distinct rows is a sum of \(n\) signs and is zero,
so \(n\) is even. We may therefore write
\begin{equation}\label{eq:initial-order}
 n=2^{2s}u^2,\qquad s\geq1,\quad u>0\ \text{odd}.
\end{equation}
The remaining task is to prove \(s=1\). A primitive character value of
a projection of \(h\) will force congruences between its coefficients.
The following lemma makes that passage from a value to coefficients
precise.

\subsection{Cyclotomic divisibility}\label{sec:local}

The coefficient rings in this lemma are localizations of cyclotomic
integer rings. Its basis and ramification properties are classical;
see, for example, \cite[pp.~320--321]{Turyn1965}. We prove the local
form, including the coefficient statement needed in both later uses.

\begin{lemma}\label{lem:local-ring}
Let \(p\) be a prime, let \(\eta\) be a primitive \(L\)-th root of unity
with \(p\nmid L\), and let \(A\) be the localization of
\(A_0=\Z[\eta]\) at a maximal ideal containing \(p\). The case \(L=1\)
is allowed. Then \(A\) has maximal ideal \(pA\), and every nonzero
element of \(A\) is a unit times a nonnegative integral power of \(p\).

For \(k\geq1\), let \(\xi\) be a primitive \(p^k\)-th root, and put
\[
 \Phi_{p^k}(X)=\sum_{j=0}^{p-1}X^{jp^{k-1}},
 \qquad D=(p-1)p^{k-1}.
\]
The ring \(A[\xi]\) is free over \(A\) with basis
\(1,\xi,\ldots,\xi^{D-1}\). It is local with maximal ideal
\((\xi-1)\), and \(p\) is a unit times \((\xi-1)^D\).
Every nonzero element is a unit times a power of \(\xi-1\).
Consequently its fraction field has an additive valuation \(v\),
normalized by
\[
 v(p)=1,\qquad v(\xi-1)=1/D.
\]
For every integer \(\ell\geq0\) and \(z\in A[\xi]\),
\begin{equation}\label{eq:valuation-ideal}
 v(z)\geq\ell\quad\Longleftrightarrow\quad z\in p^\ell A[\xi],
\end{equation}
where \(v(0)=+\infty\). In particular, for
\(R(X)=\sum_{j=0}^{D-1}r_jX^j\in A[X]\), these conditions for
\(z=R(\xi)\) are equivalent to \(r_j\in p^\ell A\) for every \(j\).
\end{lemma}

\begin{proof}
A ring generated over \(\Z\) by finitely many roots of unity is
finitely generated as an abelian group: products of bounded powers
of the roots span it. It is torsion-free because it is a subring of
\(\C\), and hence is free of finite rank. Every ideal, as a subgroup
of this free abelian group, is finitely generated. Such a ring is
therefore Noetherian, as are its localizations. This applies to \(A\)
and to \(A[\xi]\), which is a localization of \(\Z[\eta,\xi]\).

We first determine the maximal ideal of \(A\). Since \(p\nmid L\),
the polynomial \(X^L-1\) has no repeated factor over \(\F_p\).
Thus \(\F_p[X]/(X^L-1)\) is a product of fields. Its quotient
\(A_0/pA_0\) is also a product of fields. Localizing at the chosen
maximal ideal selects one field factor, so \(A/pA\) is a field and
the maximal ideal of \(A\) is \(pA\).

In a Noetherian local domain whose nonzero maximal ideal is generated
by \(t\), a nonzero nonunit can be divided by \(t\). This division
cannot continue indefinitely: the successive quotients would generate
a strictly increasing chain of principal ideals, since equality at
one step would, by cancellation, make \(t\) a unit. Thus every
nonzero element is a unit times a power of \(t\). Applied to \(A\),
this gives its integral valuation, with value \(1\) at \(p\).

We now adjoin \(\xi\). The polynomial \(\Phi_{p^k}(1+Y)\) is monic
of degree \(D\), has constant term \(p\), and reduces to \(Y^D\)
modulo \(p\). The last assertion follows by reducing
\[
 \Phi_{p^k}(X)(X^{p^{k-1}}-1)=X^{p^k}-1
\]
in characteristic \(p\). It is irreducible over \(\Frac(A)\) by
Eisenstein's argument. Here is the needed form of that argument.
For a nonzero polynomial over \(\Frac(A)\), the minimum valuation of
its coefficients is additive under multiplication: after scaling both
polynomials to minimum zero, their nonzero reductions have nonzero
product over the field \(A/pA\). If a monic integral polynomial
factors into monic polynomials, the two minima are at most zero and
sum to zero; hence both factors have coefficients in \(A\).
A nontrivial factorization of \(\Phi_{p^k}(1+Y)\) would therefore
reduce to two monic positive-degree factors of \(Y^D\). Both constant
terms would be divisible by \(p\), contradicting the constant term
\(p\) of their product.

It follows that \(A[\xi]\simeq A[X]/(\Phi_{p^k}(X))\), giving the
asserted basis. Every maximal ideal of this finite \(A\)-algebra
contracts to a maximal ideal of \(A\). Indeed, the quotient by
such a maximal ideal is a field finite as a module over the corresponding
quotient domain of \(A\). The adjugate-matrix argument for
multiplication on a finite generating set makes every element of that
field integral over the subdomain. Applying a monic relation to the
inverse of a nonzero element of the subdomain puts that inverse in the
subdomain. The subdomain is therefore a field.

Write \(\pi=\xi-1\) and \(k_0=A/pA\). The reduction
\[
 A[\xi]/pA[\xi]\simeq k_0[Y]/(Y^D)
\]
has one maximal ideal. Thus \(A[\xi]\) is local with maximal ideal
\((p,\pi)\). The equation \(\Phi_{p^k}(1+\pi)=0\) has the form
\[
 \pi^D=-p\,u_\pi,\qquad
 u_\pi=1+\sum_{j=1}^{D-1}c_j\pi^j,\qquad c_j\in A,
\]
with the sum empty when \(D=1\). The element \(u_\pi\) has residue
one and is a unit. Hence the maximal ideal is \((\pi)\), and \(p\)
is a unit times \(\pi^D\). The principal-ideal argument above now
shows that every nonzero element of \(A[\xi]\) is a unit times a power
of \(\pi\), giving the stated valuation. Since \(p^\ell\) is a unit
times \(\pi^{\ell D}\), this description proves
\eqref{eq:valuation-ideal}. The coefficient assertion follows from
uniqueness in the free \(A\)-basis.
\end{proof}
